% !TeX program = tectonic
\documentclass[11pt,a4paper]{article}
\usepackage{fontspec}
\usepackage[russian]{babel}
\babelfont{rm}[Language=Default]{Liberation Serif}
\babelfont[Language=Default]{sf}{Carlito}
\babelfont[Language=Default]{tt}{DejaVu Sans Mono}
\usepackage{amsmath,amssymb,amsthm}
\usepackage{geometry}
\geometry{a4paper, top=2.4cm, bottom=2.4cm, left=2.4cm, right=2.4cm}
\usepackage{booktabs,tabularx,enumitem,xcolor,tcolorbox}
\tcbuselibrary{breakable,skins}
\definecolor{linkblue}{HTML}{1F4E79}
\definecolor{statgreen}{HTML}{2F6B3A}
\definecolor{goldacc}{HTML}{8B7E5A}
\usepackage[colorlinks=true,linkcolor=linkblue,urlcolor=linkblue,unicode=true]{hyperref}
\theoremstyle{plain}
\newtheorem{theorem}{Теорема}
\newtheorem{lemma}[theorem]{Лемма}
\newtheorem{proposition}[theorem]{Предложение}
\newtheorem{corollary}[theorem]{Следствие}
\theoremstyle{definition}
\newtheorem{definition}[theorem]{Определение}
\theoremstyle{remark}
\newtheorem{remark}[theorem]{Замечание}
\newtcolorbox{statusbox}[1][]{colback=statgreen!5,colframe=statgreen!75,
  title={\textbf{Сводка доказанного:} #1},fonttitle=\small,boxrule=0.7pt,arc=2pt,breakable}
\newtcolorbox{databox}[1][]{colback=goldacc!6,colframe=goldacc!80,
  title={\textbf{Данные протокола:} #1},fonttitle=\small,boxrule=0.7pt,arc=2pt,breakable}
\newtcolorbox{verifybox}[1][]{colback=linkblue!5,colframe=linkblue!80,
  title={\textbf{Верификация:} #1},fonttitle=\small,boxrule=0.7pt,arc=2pt,breakable}
\widowpenalty=10000
\clubpenalty=10000
\setlength{\parskip}{0.35em}
\newcommand{\CC}{\mathbb{C}}
\newcommand{\RR}{\mathbb{R}}
\newcommand{\ZZ}{\mathbb{Z}}
\newcommand{\QQ}{\mathbb{Q}}
\newcommand{\Ga}{\Gamma}
\newcommand{\Bfun}{\mathrm{B}}
\newcommand{\im}{\operatorname{Im}}
\newcommand{\re}{\operatorname{Re}}
\newcommand{\lcm}{\operatorname{lcm}}
\newcommand{\SNF}{\operatorname{SNF}}
\newcommand{\Jac}{\operatorname{Jac}}
\newcommand{\rank}{\operatorname{rank}}
\newcommand{\Ind}{\operatorname{Index}}
\newcommand{\disc}{\operatorname{disc}}
\begin{document}
\begin{center}
{\LARGE\bfseries Теорема 2: Γ-нормировка и дробные доли π}\\[0.4em]
{\large алгебраичность нормированных периодов и радикальные константы}\\[0.6em]
{\normalsize Исаев Исхак Хамзатович}\\[0.2em]
{\small Программа <<Динамический принцип>> --- лаборатория \texttt{hodge-laboratory}}\\[0.15em]
{\small Отдельная монография теоремы · Монография 5/16}
\end{center}
\vspace{0.4em}
{\color{linkblue}\hrule height 0.8pt}
\vspace{0.8em}

\section{Постановка}

$\Ga$-нормировка поглощает трансцендентный множитель каждого
периода, превращая нормированные периоды в алгебраические числа
фазовой решётки $\frac1N\ZZ[\zeta_N]$. Требование программы ---
дробные доли $\pi$ вместо нормировок с $4\pi^2$ --- на стенде
выполняется структурно.

\begin{definition}[$\Ga$-нормировка]
Для периода $P\in\frac{\Omega_{a,b}}{N}\ZZ[\zeta_N]$ нормированный
период: $\widehat{P}=P/\Omega_{a,b}\in\frac1N\ZZ[\zeta_N]$.
\end{definition}

\begin{theorem}[$\Ga$-нормировка и дробные доли $\pi$]\label{th:norm}
Пусть $N\in\{15,30\}$. Тогда:
(1) алгебраичность: $\widehat P\in\frac1N\ZZ[\zeta_N]$ --- фазовая
решётка имеет знаменатель $N$;
(2) редукция: каждое $\Omega_{a,b}$ выражается через значения $\Ga$
в канонических позициях и множители $\pi/\sin(\pi k/N)$; при
$a+b=N$ --- чистая доля $\pi$;
(3) отсутствие $4\pi^2$: в замкнутой форме $\pi$ входит только через
$\pi/\sin(\pi k/N)$ и формулу умножения; никакая степень $4\pi^2$ не
возникает в нормировке;
(4) радикальные константы:
\begin{align}
  b_{Ch}(15)&=1-\cos\frac{2\pi}{15}
  =1-\frac{1+\sqrt5+\sqrt3\,\sqrt{10-2\sqrt5}}{8},\\
  b_{Ch}(30)&=1-\cos\frac{\pi}{15}
  =1-\frac{\sqrt{30+6\sqrt5}+\sqrt5-1}{8}.
\end{align}
\end{theorem}

\begin{proof}
(1) --- следствие замкнутой формы периода: деление решётки
$\frac{\Omega}N\ZZ[\zeta_N]$ на $\Omega$.
(2)--(3) --- инвентарь соотношений: отражение вводит $\pi$ только
как $\pi/\sin(\pi k/N)$, умножение --- как $(2\pi)^{(m-1)/2}$ при
целых $m$; нормировка строится из самих периодов, а не спускается
из L-функций, поэтому $4\pi^2$ не порождается. (4):
$\cos24^\circ=\cos(60^\circ-36^\circ)$ с
$\cos36^\circ=\frac{1+\sqrt5}4$,
$\sin36^\circ=\frac{\sqrt{10-2\sqrt5}}4$ даёт формулу для $15$;
$\cos12^\circ=\cos(30^\circ-18^\circ)$ с
$\cos18^\circ=\frac{\sqrt{10+2\sqrt5}}4$,
$\sin18^\circ=\frac{\sqrt5-1}4$ --- для $30$.
Численно: $b_{Ch}(15)\approx0.086454$, $b_{Ch}(30)\approx0.021854$.
\end{proof}

\begin{remark}
Доля $\pi$ всегда рациональна со знаменателем $N$: множитель
$\pi/\sin(\pi k/N)$ несёт угол $k\pi/N$; фазовая решётка несёт
знаменатель $N$ в алгебре. Структурное утверждение --- верно для
всех блоков сразу. Лестница $\pi/7$, $\pi/15$, $\pi/30$ --- ступени
одной лестницы $\mu_N$-квантов (сертификаты G, H).
\end{remark}

\begin{databox}[числа]
$b_{Ch}(7)\approx0.7526$; $b_{Ch}(9)\approx0.1172$;
$b_{Ch}(11)\approx0.0823$; $b_{Ch}(15)\approx0.086454$;
$b_{Ch}(30)\approx0.021854$. Универсальный Tate-фактор $4\pi^2$
живёт в спектре решёток (G1), не в нормировке периодов.
\end{databox}

\begin{statusbox}[итог]
Теорема фиксирует: знаменатель фазовой решётки равен уровню $N$;
трансцендентная часть редуцируется к дробным долям $\pi$;
радикальные константы стенда выписаны в замкнутой форме.
\end{statusbox}

\begin{verifybox}[как воспроизвести]
\texttt{python3 laboratory.py} $\to$ <<Стенд N=15/30>> (теорема 2
блока); Lean: \texttt{bch\_radicals} --- сверка радикальных
выражений; диаграмма \texttt{fig\_bch.png}.
\end{verifybox}

\vfill
{\color{linkblue}\hrule height 0.6pt}
\vspace{0.5em}
{\small\color{gray}
Лицензия: индивидуальная исключительная лицензия Исаева Исхака Хамзатовича.
Все права на вычисления, формулы и тексты принадлежат автору.
Верификация: \texttt{python3 laboratory.py} (пункт <<Стенды>>/<<Сертификаты>>),
Lean: \texttt{verification/lean/HodgeLaboratory.lean}.
}
\end{document}
